LLM judges are often anonymised on the assumption that they cannot tell who
wrote a text. The usual check is whether a judge names itself as the author
of its own text more often than chance. We show that this check is
unreliable. A high self-naming rate can mean the judge recognises its
writing, but it can also mean the text is easy for anyone to attribute, or
that the judge answers ``me'' to almost everything. We therefore compare
each judge with two references from the same experiment: how often it names
itself on text it did not write, and how often the \emph{other} judges name
it on its text. Across nine judges (five frontier, four open-weight), three
sets of texts and three question formats, the usual check finds
self-recognition in \SVPStd{} of \SVPCells{} cases; only \SVPProper{} hold up
against both references, all of them Claude or GPT on one set of prose
texts, and none on a larger set of 600 texts in four genres. Asked directly
whether they wrote a text, three frontier judges deny writing any of them,
their own included. Self-recognition should be measured for each judge, on
the texts it will actually judge, against a peer baseline.
